Mechanistic studies of large language models localize function in parameters, heads or neurons, but those parts do not specify how one inference unfolds. We treat the ordered layer-state trajectory as the empirical process. Across representative Qwen, Llama and Gemma checkpoints, realized order was 1.60--2.24 times more chord-aligned than endpoint-preserving shuffles. This geometry recurred in a four-choice benchmark and two controlled task families, but was stronger in three random initializations of each architecture, revealing an architecture-compatible scaffold rather than learned reasoning by itself. Ordered partial trajectories forecast a held-out candidate margin ($R^2=0.39$--$0.71$; AUROC $=0.84$--$0.98$). After additive interventions failed, confidence-gated local-transition actuators crossed a declared full-vocabulary boundary in 17/87, 13/87 and 9/87 held-out prompts, with 0/203, 2/203 and 0/203 non-target changes. Layer order thus defines an observable, computable and controllable inference process, while task coordinates, action assignment and answer correctness remain bounded.
